% Kevin integra; cada bloque aporta sus figuras. ~1 pagina.
% Las dos tablas las genera proyecto/pc/campana.py: no editarlas a mano aqui.
\section{Resultados}
\label{sec:resultados}

\subsection{Criterios de aceptación}
Los criterios se fijaron antes de medir y quedaron versionados en
\texttt{config/pruebas.json}; no se ajustaron al ver los resultados. El piso
general es $\Delta\mathrm{SNR}\geq3$~dB, porque una mejora menor no se distingue
a oído ni sobresale de la variabilidad entre repeticiones, y conservar al menos
el 80\,\% de la energía útil, medida en $\pm60$~Hz de cada componente de la
señal. Cada prueba endurece o relaja ese piso según cuánto se solapen señal y
ruido (Tabla~\ref{tab:criterios}): P1 exige 15~dB porque la interferencia está
separada de la banda útil, y P5 exige solamente no empeorar, porque ahí señal y
ruido comparten frecuencia y el límite lo impone el problema, no el método. La
latencia extremo a extremo se acota a 150~ms, que deja margen sobre los 96~ms de
ventana y salto de un bloque de 1024 muestras a 16~kHz.

\begin{table}[t]
\caption{Criterios de aceptación fijados antes de medir.}
\label{tab:criterios}
\centering\small
\IfFileExists{../resultados/tabla_criterios.tex}
  {% generada por proyecto/pc/campana.py a partir de config/pruebas.json
\begin{tabular}{llrrr}
\toprule
Prueba & Condición & $\Delta$SNR mín. & $E_{\text{útil}}$ mín. & Aten. mín. \\
 & & (dB) & (\%) & (dB) \\
\midrule
P1 & interferencia sinusoidal separada & 15 & 80 & 30 \\
P2 & ruido de alta frecuencia & 10 & 80 & --- \\
P3 & ruido de banda ancha & 3 & 80 & --- \\
P4 & ruido impulsivo & 3 & 80 & --- \\
P5 & superposición parcial señal-ruido & 0 & 70 & --- \\
P6 & definida por el profesor en la defensa & 3 & 80 & --- \\
\bottomrule
\end{tabular}
}
  {\pendiente{correr \texttt{pc/campana.py} para generar la tabla}}
\end{table}

\subsection{Enlace entre nodos}
Cada trama añade nueve bytes de control (sincronización, tipo, secuencia,
longitud y CRC-16/CCITT) a una carga de 2048~bytes, un 0,44\,\% de sobrecarga.
A 16~kHz y 921\,600~baud el flujo de muestras ocupa cerca del 35\,\% del enlace,
de modo que la numeración de secuencia y la detección de errores no compiten con
el ancho de banda útil. Las tramas de estado se envían a lo sumo una vez por
segundo y solo cuando algún contador cambia, con un costo de 21~bytes por
segundo. El valor de esto no es el ahorro: la secuencia permite contar bloques
perdidos y el CRC descartar bloques corruptos, así que una métrica degradada se
atribuye al filtrado y no a una pérdida de datos silenciosa.

\subsection{Campaña experimental}
Cada caso se repitió tres veces y se reporta el promedio con su desviación
estándar (Tabla~\ref{tab:campana}). Las métricas se calculan en la computadora
sobre lo que registró \texttt{pc/registrar.py}, con la misma implementación de
MSE, SNR y energía por bandas usada en la caracterización previa, y después de
compensar por correlación cruzada el retardo entre la referencia y lo capturado.
La ganancia de la cadena PWM--RC--ADC es desconocida, así que se estima por
mínimos cuadrados dentro de la banda útil; estimarla sobre todo el ancho de banda
la contamina con la interferencia y descalibra la referencia.

\begin{table}[t]
\caption{Resultados de la campaña: promedio $\pm$ desviación estándar de tres
repeticiones.}
\label{tab:campana}
\centering\small
\IfFileExists{../resultados/tabla_campana.tex}
  {% generada por proyecto/pc/campana.py
\begin{tabular}{llrrrrc}
\toprule
Prueba & Método & $\Delta$SNR (dB) & MSE & $E_{\text{útil}}$ (\%) & Aten. (dB) & Cumple \\
\midrule
P3 & temporal & $4.18 \pm 0.32$ & $1.63e-04$ & 70.7 & 42.8 & no \\
P3 & espectral & $5.22 \pm 0.10$ & $1.29e-04$ & 54.0 & 6.4 & no \\
\bottomrule
\end{tabular}
}
  {\pendiente{campaña sin correr: falta la evidencia de P1--P5}}
\end{table}

\subsection{Comparación de los dos métodos}
La caracterización de escritorio que precedió al diseño ya anticipaba que ninguno
de los dos métodos gana en todos los escenarios, y es la razón por la que el
sistema conmuta entre ambos en vez de elegir uno. Con una interferencia tonal
fuera de la banda de paso, el filtrado temporal mejoró la SNR en 27,5~dB frente
a 22,2~dB del espectral; con ruido blanco y con solapamiento parcial, en cambio,
el temporal no aportó mejora apreciable (0,18 y 0,00~dB) mientras el espectral
mantuvo 5,5 y 0,5~dB. Esa uniformidad del método espectral se paga: en el caso
de ruido blanco costó un 13,3\,\% de la energía útil, mientras el temporal
conserva casi toda la que deja pasar su banda. El compromiso
ruido--distorsión es, entonces, asimétrico entre los dos métodos, y el
diagnóstico del Nodo~1 es lo que decide cuál de los dos conviene por bloque.

\pendiente{Kevin: latencia medida, memoria y carga de cada nodo, con los datos
que reporten los bloques 1, 2 y 3 al integrar en placa.}
